\begin{figure}[H]
    \centering
    {
    \chemnameinit{\chemfig{CH_3-C(-[6]H)(-[2]{\color{red}H})-O{\color{red}H}}}
    \schemestart
        \chemname{\chemfig{CH_3-C(-[6]H)(-[2]{\color{red}H})-O{\color{red}H}}}{\small\bfseries Etanol}
        \arrow(.base east--.base west){-U>[NAD+][NAD{\color{red}H} + {\color{red}H+}][][0.35][60]}[,2.5] 
        \chemname{\chemfig{CH_3-C(-[7]H)=[1]O}}{\small\bfseries Acetaldehido}
    \schemestop
    \chemnameinit{}
    \vspace*{1em} \par
    }
    \rule{\linewidth}{0.2pt}
    \caption[Oxidación del etanol en acetaldehido por la alcohol deshidrogenasa]{Reacción catalizada por la alcohol deshidrogenasa (EC 1.1.1.). Aquí la enzima utiliza el NAD\textsuperscript{+} como cofactor para oxidar el etanol, esto al eliminar de este 2 átomos de hidrógeno y 2 electrones, lo que transforma al etanol en acetaldehido. La alcohol deshidrogenasa está presente en nuestro hígado, y hoy día el principal uso que hacemos de ella es para el metabolismo de bebidas alcohólicas.}
    \label{fig:reaccion_alcohol_deshidrogenasa}
\end{figure}